%% GENERATED by tools/from_docs.py from stepss-docs. Do not edit by hand:
%% edit the page named below in stepss-docs and regenerate.
%% source: stepss-docs/src/content/docs/models/discrete-controllers/monitoring.md

\chapter{Monitoring controllers}\label{doc:model-dctl-mon}

The \texttt{DCTL} record these models are entered with, and the rules on model names, are on the \hyperref[doc:model-dctl]{Discrete Controller Models} page.

\section{\texorpdfstring{SIM\_MINMAXSPEED (\texttt{dctl\_sim\_minmaxspeed}): Speed Monitoring (Min/Max)}{SIM\_MINMAXSPEED (dctl\_sim\_minmaxspeed): Speed Monitoring (Min/Max)}}\label{model-dctl-mon:sim_minmaxspeed-dctl_sim_minmaxspeed-speed-monitoring-minmax}

\subsection{Description}\label{model-dctl-mon:description}

A \textbf{simulation monitoring} controller that checks whether any synchronous machine speed violates upper or lower bounds. If a violation persists beyond a deadtime, the controller logs a warning or stops the simulation, depending on the \texttt{stop\_sim} flag. It is used to automatically terminate simulations that have become numerically unstable or that have reached an unacceptable operating state.

\subsection{Logic Description}\label{model-dctl-mon:logic-description}

At each evaluation, the controller scans all synchronous machines and checks:

\[\text{if } \omega > \omega_{max} \text{ or } \omega < \omega_{min}: \quad \text{record violation time}\]
\[\text{if violation duration} > T_{deadtime}: \quad \text{warn or stop simulation}\]

\subsection{Parameters}\label{model-dctl-mon:parameters}

\begin{small}\begin{longtable}[]{@{}
  >{\raggedright\arraybackslash}p{(\linewidth - 6\tabcolsep) * \real{0.2690}}
  >{\raggedright\arraybackslash}p{(\linewidth - 6\tabcolsep) * \real{0.2437}}
  >{\raggedright\arraybackslash}p{(\linewidth - 6\tabcolsep) * \real{0.4873}}
@{}}
\toprule\noalign{}
\begin{minipage}[b]{\linewidth}\raggedright
Field
\end{minipage} & \begin{minipage}[b]{\linewidth}\raggedright
Working variable
\end{minipage} & \begin{minipage}[b]{\linewidth}\raggedright
Description
\end{minipage} \\
\midrule\noalign{}
\endhead
\bottomrule\noalign{}
\endlastfoot
\texttt{omega\_max} & \texttt{w(1)} & Upper speed limit (pu) \\
\texttt{omega\_min} & \texttt{w(2)} & Lower speed limit (pu) \\
\texttt{deadtime} & \texttt{w(3)} & Deadtime before action (s) \\
\texttt{stop\_sim} & \texttt{w(4)} & 1 = stop simulation on violation, 0 = warn only \\
\end{longtable}\end{small}

4 fields required. Internal: \texttt{w(5)} = last violation time.

Parameter names exported: \texttt{SPEEDMIN}, \texttt{SPEEDMAX}, \texttt{DEADTIME}.

\subsection{Usage Example}\label{model-dctl-mon:usage-example}

\begin{Verbatim}
DCTL  SIM_MINMAXSPEED  SPEEDMON
      # SPEEDMIN: lowest machine speed allowed (pu)
      0.5
      # SPEEDMAX: highest machine speed allowed (pu)
      1.2
      # start_time: monitoring starts at this time (s)
      0.5
      # warn_only: 0 = stop the simulation on a violation, otherwise log a warning only
      0
;
\end{Verbatim}

\begin{center}\rule{0.5\linewidth}{0.5pt}\end{center}

\section{\texorpdfstring{SIM\_MINMAXVOLT (\texttt{dctl\_sim\_minmaxvolt}): Voltage Monitoring (Min/Max)}{SIM\_MINMAXVOLT (dctl\_sim\_minmaxvolt): Voltage Monitoring (Min/Max)}}\label{model-dctl-mon:sim_minmaxvolt-dctl_sim_minmaxvolt-voltage-monitoring-minmax}

\subsection{Description}\label{model-dctl-mon:description-1}

Analogous to \texttt{dctl\_sim\_minmaxspeed} but for \textbf{bus voltages}. The controller scans all buses and checks whether any voltage violates the specified bounds. If the violation persists for longer than the deadtime, a warning is logged or the simulation is terminated.

\subsection{Logic Description}\label{model-dctl-mon:logic-description-1}

\[\text{if } V > V_{max} \text{ or } V < V_{min}: \quad \text{record violation time}\]
\[\text{if violation duration} > T_{deadtime}: \quad \text{warn or stop simulation}\]

\subsection{Parameters}\label{model-dctl-mon:parameters-1}

\begin{small}\begin{longtable}[]{@{}lll@{}}
\toprule\noalign{}
Field & Working variable & Description \\
\midrule\noalign{}
\endhead
\bottomrule\noalign{}
\endlastfoot
\texttt{V\_max} & \texttt{w(1)} & Upper voltage limit (pu) \\
\texttt{V\_min} & \texttt{w(2)} & Lower voltage limit (pu) \\
\texttt{deadtime} & \texttt{w(3)} & Deadtime before action (s) \\
\texttt{stop\_sim} & \texttt{w(4)} & 1 = stop on violation, 0 = warn only \\
\end{longtable}\end{small}

4 fields required. Internal: \texttt{w(5)} = last violation time.

Parameter names exported: \texttt{VMIN}, \texttt{VMAX}, \texttt{DEADTIME}.

\subsection{Usage Example}\label{model-dctl-mon:usage-example-1}

\begin{Verbatim}
DCTL  SIM_MINMAXVOLT  VOLTMON
      # VMIN: lowest bus voltage allowed (pu)
      0.80
      # VMAX: highest bus voltage allowed (pu)
      1.15
      # start_time: monitoring starts at this time (s)
      1.0
      # warn_only: 0 = stop the simulation on a violation, otherwise log a warning only
      0
;
\end{Verbatim}

\begin{center}\rule{0.5\linewidth}{0.5pt}\end{center}

\section{\texorpdfstring{VOLT\_VAR (\texttt{dctl\_voltage\_variability}): Voltage Variability Monitor}{VOLT\_VAR (dctl\_voltage\_variability): Voltage Variability Monitor}}\label{model-dctl-mon:volt_var-dctl_voltage_variability-voltage-variability-monitor}

The data-file model name is \texttt{VOLT\_VAR}.

\subsection{Description}\label{model-dctl-mon:description-2}

Computes a \textbf{running Exponential Moving Average (EMA)} of voltage variability (variance) across a list of monitored buses. It is used to assess power quality metrics such as voltage flicker or ramping caused by variable renewable energy sources.

The EMA uses two decay constants (\(\lambda_1\), \(\lambda_2\)) derived from the averaging time window, enabling a computationally efficient recursive estimation:

\[\text{EMA}(t) = \lambda_1 \cdot \text{EMA}(t-1) + (1 - \lambda_1) \cdot V(t) + (\lambda_1 - \lambda_2)(V(t) - V(t-1))\]

The controller accumulates per-bus variances and a system-wide sum.

\subsection{Parameters}\label{model-dctl-mon:parameters-2}

\begin{small}\begin{longtable}[]{@{}
  >{\raggedright\arraybackslash}p{(\linewidth - 4\tabcolsep) * \real{0.3129}}
  >{\raggedright\arraybackslash}p{(\linewidth - 4\tabcolsep) * \real{0.6871}}
@{}}
\toprule\noalign{}
\begin{minipage}[b]{\linewidth}\raggedright
Field
\end{minipage} & \begin{minipage}[b]{\linewidth}\raggedright
Description
\end{minipage} \\
\midrule\noalign{}
\endhead
\bottomrule\noalign{}
\endlastfoot
\texttt{aver\_time\_window} & Averaging time window (s) for EMA computation \\
\texttt{nb\_list} & Number of buses to monitor (bus names follow in subsequent RAMSES data records) \\
\end{longtable}\end{small}

2 fields required. Only one \texttt{VOLT\_VAR} controller may exist per simulation. Parameters are stored in module-level variables (\texttt{volt\_var\_mod}) rather than \texttt{w(*)}.

\subsection{Usage Example}\label{model-dctl-mon:usage-example-2}

\begin{Verbatim}
DCTL  VOLT_VAR  VVAR
      # window: averaging time window (s)
      60.0
      # nb_list: number of most variable buses written to the log at the end (an integer)
      5
;
\end{Verbatim}
